\documentclass[review,3p]{elsarticle}

%% ============================================================
%% PAPER 1: TOPOLOGY ABSTRACTION EFFECTS
%% Target Journal: Energy Conversion and Management
%% Version: 4.0 — Restructured for clarity and reduced redundancy
%% ============================================================

\usepackage[utf8]{inputenc}
\usepackage[T1]{fontenc}
\usepackage{amsmath,amssymb,amsfonts}
\usepackage{graphicx}
\usepackage{booktabs}
\usepackage{multirow}
\usepackage{array}
\usepackage{siunitx}
\usepackage{xcolor}
\usepackage{hyperref}
\usepackage{cleveref}
\usepackage{lineno}
\usepackage{float}
\usepackage{subfig}
\usepackage{enumitem}
\usepackage[version=4]{mhchem}

\modulolinenumbers[5]
\sisetup{per-mode=symbol}
\hypersetup{colorlinks=true, linkcolor=blue, citecolor=blue, urlcolor=blue}

\newcommand{\placeholder}[1]{\textcolor{red}{[#1]}}

\journal{Energy Conversion and Management}

\begin{document}

\begin{frontmatter}

\title{Topology Abstraction Effects on Dispatch Optimization of
Electrified District Heating Networks: A Controlled MILP Comparison}

\author[tum]{Lukas [Nachname]\corref{cor1}}
\ead{lukas.nachname@tum.de}
\author[tum]{[Betreuer]}
\author[partner]{[Industriepartner]}

\cortext[cor1]{Corresponding author.}

\address[tum]{Lehrstuhl für Energiewirtschaft und Anwendungstechnik,
  Technische Universität München, München, Germany}

\address[partner]{[Industriepartner Institution], [Stadt], Austria}

%% ============================================================
%% ABSTRACT
%% ============================================================
\begin{abstract}
District heating network planning studies require practitioners to choose
a topology representation level, yet no quantitative guidance exists for
this decision. Prior comparisons confound topology abstraction with
physics model changes, rendering it impossible to attribute observed
cost differences to either factor independently.

This paper presents the first controlled comparison that isolates
topology as the sole experimental variable. Three abstraction levels
are defined within a unified mixed-integer linear programming (MILP)
framework: L1 (copperplate, no network physics), L2 (simplified
multi-node with physics-based thermal losses), and L3 (detailed
multi-node with identical physics to L2). The framework is validated
against measured operational data from an Austrian industrial heat
network comprising heat pumps, thermal storage, a gas boiler, an
electrode boiler, and grid electricity coupling. Generalizability is
assessed across 36 synthetic network configurations spanning systematic
variations in pipe length, demand heterogeneity, node count, and
storage capacity.

Results quantify the topology effect: the L1$\rightarrow$L2 transition
(physics inclusion) drives cost differences of \placeholder{X}\%
(\placeholder{€Y,000}/year) and \ce{CO2} differences of
\placeholder{Z}~t/year, while the L2$\rightarrow$L3 transition
(spatial refinement) adds \placeholder{$<$X}\% for uniform-demand
networks but up to \placeholder{Y}\% for spatially heterogeneous
systems. Thermal storage dispatch emerges as the primary diagnostic
variable for model adequacy: L1 systematically misrepresents charging
patterns by ignoring loss-dependent temporal holding costs, shifting
optimal charge timing by \placeholder{X}~hours on average in winter.
The synthetic analysis reveals that the L1$\rightarrow$L2 gap scales
with total pipe length ($R^2\!=\,$\placeholder{XX}), while the
L2$\rightarrow$L3 gap is driven by demand heterogeneity. These
findings yield actionable selection criteria: L2 is the minimum
fidelity for storage-containing systems, while L3 is warranted when
the demand heterogeneity index exceeds \placeholder{0.X}.
\end{abstract}

\begin{keyword}
District heating \sep Mixed-integer linear programming \sep
Topology abstraction \sep Thermal storage dispatch \sep
Heat pump optimization \sep Network modeling fidelity
\end{keyword}

\end{frontmatter}

\linenumbers

%% ============================================================
%% 1. INTRODUCTION
%% ============================================================
\section{Introduction}
\label{sec:introduction}

\subsection{Motivation and Problem Statement}
\label{sec:motivation}

Space and process heat account for approximately 50\% of final energy
consumption in the European Union, of which roughly 60\% is still
supplied by fossil fuels~\cite{EUheating2016}. District heating (DH)
networks offer a systemic decarbonization pathway through sector
coupling: large-scale heat pumps, combined heat and power (CHP) units,
thermal energy storage, and electrode boilers can be coordinated within
a single network~\cite{Lund2014_4GDH,Werner2017}. Planning and
operational optimization of such multi-source networks increasingly
relies on mixed-integer linear programming (MILP) due to its global
optimality guarantees and tractable handling of binary commitment
decisions.

Every MILP-based planning study, however, confronts a fundamental
modeling choice that has received remarkably little systematic
attention: the level of spatial detail at which the network is
represented. At one extreme, \emph{copperplate} models aggregate all
components and demands to a single bus, ignoring spatial structure and
thermal losses entirely. At the other extreme, \emph{full-topology}
models resolve the physical pipe-by-pipe network, capturing thermal
losses that typically constitute 5--15\% of annual heat supply in
Central European networks~\cite{Frederiksen2013,Talebi2016}. The
practical consequences of this choice affect cost estimation, dispatch
quality, and \ce{CO2} accounting simultaneously---yet practitioners
currently lack quantitative guidance on the accuracy--tractability
trade-off.

\subsection{Research Gap}
\label{sec:gap}

The core problem is \emph{confounding}: all published comparisons
between different topology representations change the physics model
and the spatial aggregation level simultaneously, making it impossible
to attribute observed differences to either factor independently. For
instance, comparing the full nonlinear pipe model with
temperature-dependent COP used by Vesterlund et al.~\cite{Vesterlund2017}
against the copperplate formulation with constant thermal efficiencies
of Gabrielli et al.~\cite{Gabrielli2018} conflates topology effects
with (1)~NLP vs.\ constant-loss formulations, (2)~temperature-dependent
vs.\ constant COP, and (3)~potentially differing temporal resolutions.
The same pattern appears systematically across the literature
(\Cref{tab:prior_work} in \Cref{sec:related_work}). As a result, the
practitioner question that matters most---\emph{how much does topology
abstraction itself change cost and emission estimates, for a system
with fixed physics?}---remains unanswered.

Three specific gaps persist:
\begin{enumerate}[nosep]
    \item No study isolates topology abstraction as the sole
    experimental variable while holding physics constant.
    \item No quantitative guidance exists for topology selection as a
    function of observable network properties.
    \item The mechanism by which topology level affects storage
    dispatch---through loss-dependent temporal holding costs---has not
    been analyzed.
\end{enumerate}

\subsection{Contributions and Research Questions}
\label{sec:contributions}

This paper addresses these gaps through four contributions:

\begin{enumerate}
    \item \textbf{L1/L2/L3 taxonomy}: A reusable classification
    framework for DH modeling fidelity---copperplate (L1), simplified
    multi-node (L2), multi-node MILP (L3) and multi-node MIQP (L3) ---enabling
    consistent comparison across future studies (\Cref{sec:topology_def}).

    \item \textbf{First controlled comparison}: Three topology levels
    embedded within a unified MILP framework where L2 and L3 share
    identical physics, validated against measured industrial operational
    data (\Cref{sec:results}).

    \item \textbf{Storage dispatch as primary diagnostic}: Mechanistic
    explanation of how topology affects cost via loss-dependent temporal
    holding costs, with storage charge timing as a directly observable
    adequacy indicator (\Cref{sec:results_storage}).

    \item \textbf{Actionable decision criteria}: Parametric analysis
    across 36 synthetic configurations yielding topology-selection
    rules based on observable network properties
    (\Cref{sec:results_generalizability}).
\end{enumerate}

These contributions answer three research questions:
\begin{description}[style=nextline,nosep]
    \item[RQ1 --- Cost Impact:]
    How much does topology abstraction change annual operational cost
    and \ce{CO2} emissions when physics are held constant?

    \item[RQ2 --- Dispatch Sensitivity:]
    Which dispatch decisions are most sensitive to topology level,
    and what is the physical mechanism?

    \item[RQ3 --- Generalizability:]
    Which network properties determine whether detailed topology
    representation is warranted?
\end{description}

%% ============================================================
%% 2. RELATED WORK
%% ============================================================
\section{Related Work}
\label{sec:related_work}

\subsection{MILP Optimization of District Heating Systems}
\label{sec:rw_milp}

MILP has become the dominant methodology for operational optimization
of district heating systems~\cite{Fang2016,Guelpa2019review}. Several
open-source frameworks support multi-carrier energy system optimization:
oemof~\cite{Hilpert2018_oemof} with the oemof-thermal
extension~\cite{Krien2020_oemofthermal},
PyPSA~\cite{Brown2018_pypsa}, FINE~\cite{Welder2018_FINE},
Calliope~\cite{Pfenninger2018_calliope}, urbs~\cite{Dorfner2016_urbs},
and Balmorel~\cite{Wiese2018_balmorel}. Despite their diversity, these
frameworks share common limitations for detailed DH analysis: thermal
losses are typically time-invariant, heat pump COP is constant or
seasonal, and network topology is either fully aggregated or fixed at
a single detail level.

\subsection{Topology Representation Spectrum}
\label{sec:rw_topology}

The DH modeling literature exhibits a continuous spectrum of spatial
aggregation, which we organize into three categories corresponding
to our L1/L2/L3 taxonomy.

\textbf{Single-node aggregation (L1-equivalent):} Standard in
national-scale energy system models where heating is one of many
sectors~\cite{Connolly2014_EnergyPLAN,Brown2018_pypsa,
Lund2017_SmartEnergy}. No spatial structure or thermal losses.

\textbf{Zone-based models (L2-equivalent):} Demand aggregated into
geographic zones (3--15 nodes) connected by simplified pipe
segments~\cite{Wirtz2020,Mertz2016,Buffa2019}. Bünning et
al.~\cite{Buenning2020} demonstrate that zone-level aggregation
captures 85--95\% of dynamic temperature variation.

\textbf{Full-graph models (L3-equivalent):} Node-by-node
representation with pipe-by-pipe parameters from GIS
data~\cite{Vesterlund2017,Schweiger2017,vanderHeijde2017,
Blommaert2020}. Giraud et al.~\cite{Giraud2015} show tractability
up to $\approx$100 nodes for hourly annual horizons.

\Cref{tab:prior_work} positions representative studies. The critical
observation is that in all prior work, the physics model changes
simultaneously with the topology level---the ``Controlled?'' column
is uniformly ``No.''

\begin{table*}[htbp]
\centering
\caption{Representative DH optimization studies by topology level
and physics model. No prior study isolates topology from physics.}
\label{tab:prior_work}
\small
\begin{tabular}{@{}p{3.5cm}p{2.0cm}p{3.0cm}p{2.2cm}p{1.5cm}p{1.8cm}@{}}
\toprule
Reference & Topology & Physics Model & COP Treatment & Scale & Controlled? \\
\midrule
Connolly et al.~\cite{Connolly2014_EnergyPLAN}
  & Copperplate & Fixed efficiency & Constant seasonal & National & No \\
Gabrielli et al.~\cite{Gabrielli2018}
  & Copperplate & Fixed efficiency & Constant & 10~MW & No \\
Mertz et al.~\cite{Mertz2016}
  & Zone (8 nodes) & Linear losses & Constant & 20~MW & No \\
Wirtz et al.~\cite{Wirtz2020}
  & Zone (5 nodes) & Linear losses & Fixed seasonal & 25~MW & No \\
Bünning et al.~\cite{Buenning2020}
  & Reduced-order & Simplified & Constant & 8~MW & No \\
Giraud et al.~\cite{Giraud2015}
  & Full graph & Steady-state & Temperature-dep. & 30~MW & No \\
Vesterlund et al.~\cite{Vesterlund2017}
  & Full graph & NLP thermal & Temperature-dep. & 50~MW & No \\
Schweiger et al.~\cite{Schweiger2017}
  & Full graph & Modelica co-sim & Dynamic & 15~MW & No \\
van der Heijde et al.~\cite{vanderHeijde2017}
  & Full graph & Pipe model comp. & -- & Review & Partial$^{a}$ \\
Blommaert et al.~\cite{Blommaert2020}
  & Full graph & Adjoint NLP & Temperature-dep. & 30~MW & No \\
\textbf{This work}
  & \textbf{L1/L2/L3} & \textbf{Identical MILP} & \textbf{Pre-computed}
  & \textbf{$\approx$12~MW} & \textbf{Yes} \\
\bottomrule
\multicolumn{6}{@{}l}{\footnotesize $^{a}$~Compared pipe loss
formulations but not spatial aggregation levels.}
\end{tabular}
\end{table*}

\subsection{Thermal Losses and Storage Dispatch}
\label{sec:rw_losses_storage}

Pipe thermal losses follow from steady-state or dynamic heat transfer
between fluid and soil. The steady-state linear model, in which loss
is proportional to the fluid--ground temperature
difference~\cite{Paalsson1999,Frederiksen2013}, forms the basis for
MILP-compatible loss calculations. Werner~\cite{Werner2017} reports
European DH loss ratios of 5--25\%, with modern pre-insulated networks
typically below 10\%. Talebi et al.~\cite{Talebi2016} confirm that
loss fraction increases with supply temperature and pipe length.

Thermal energy storage (TES) provides flexibility through peak
shaving, load shifting, and renewable
integration~\cite{Vandermeulen2018,Guelpa2019review}. Christidis et
al.~\cite{Christidis2017} demonstrate that optimal TES dispatch
shifts heat production by 4--8~hours to exploit electricity price
differentials---a pattern directly dependent on accurate temporal cost
signals. In copperplate models, the only time-varying cost is the
electricity price. In network-aware models, losses introduce an
additional implicit holding cost: heat stored for later discharge
must additionally compensate for losses that are highest during cold
high-demand periods. This mechanism has not been quantified in a
controlled setting---a gap this paper fills.
\textbf{Hypothesis:} Storage dispatch is the most sensitive 
indicator of topology adequacy because temporal arbitrage 
depends on holding costs invisible to L1. This is tested 
quantitatively in \Cref{sec:results_storage}.
%% ============================================================
%% 3. METHODOLOGY
%% ============================================================
\section{Methodology}
\label{sec:methodology}

This section presents the controlled experimental design
(\Cref{sec:exp_design}), the unified MILP formulation
(\Cref{sec:milp}), the demand heterogeneity metric
(\Cref{sec:model_hi}), and implementation details
(\Cref{sec:model_impl}). The key design principle is that topology
is the sole experimental variable: L2 and L3 share identical physics,
while L1 deliberately omits network losses as a baseline.

\subsection{Controlled Experimental Design}
\label{sec:exp_design}

\subsubsection{Design Principles}

The physics model is deliberately simplified and---critically---held
\emph{identical} across L2 and L3:
\begin{itemize}[nosep]
    \item \textbf{Steady-state thermal losses}: Per-pipe, separately
    for supply and return (\Cref{sec:model_losses}).
    \item \textbf{Fixed supply temperature}: Determined by a heating
    curve, not optimized. Eliminates temperature--flow coupling while
    preserving seasonal loss variation.
    \item \textbf{Pre-computed COP}: Time-varying parameters from
    bilinear interpolation of manufacturer data
    (\Cref{sec:model_hp}), entering as parameters to preserve MILP
    linearity.
    \item \textbf{No transport delays}: Instantaneous heat
    propagation within each hourly timestep (estimated transport
    times: 0.2--1.8~h). Delays would confound topology with temporal
    dynamics.
    \item \textbf{No pressure optimization}: Pumping power confirmed
    at $<$2\% of thermal losses via post-processing
    (Darcy--Weisbach at peak flow).
    \item \textbf{Dispatch-only}: All component capacities fixed.
\end{itemize}

These are \emph{design choices for experimental control}, not
framework limitations. The underlying framework supports extended
physics including transport delays, pressure drop coupling, and unit
commitment, documented in a companion paper~\cite{Paper2_CALION}.

\subsubsection{L1/L2/L3 Topology Definitions}
\label{sec:topology_def}

\textbf{Level~1 (Copperplate):} Single bus, no spatial structure, no
pipe losses ($Q_{\mathrm{loss}}=0$ for all $t$). A single scalar heat
balance.

\textbf{Level~2 (Simplified multi-node):} Aggregated demand zones
(5 nodes in the primary case) connected by simplified pipe segments
with physics-based steady-state losses. The total $U \cdot L$ product
is preserved from L3, ensuring matched aggregate annual loss energy.

\textbf{Level~3 (Detailed multi-node):} Full spatial resolution from
GIS data (30 nodes in the primary case). Individual pipe parameters.
Same loss formulation as L2.

\textbf{Controlled matching:} L2 and L3 use identical loss
formulations, identical total installed capacities per component
type, identical storage parameters, and identical boundary conditions.
Only node count and spatial pipe resolution differ.
\Cref{tab:feature_matrix} specifies feature activation by level.

\begin{table}[htbp]
\centering
\caption{Feature activation across topology levels.
L2 and L3 share identical formulations for all active features.}
\label{tab:feature_matrix}
\begin{tabular}{@{}lccc@{}}
\toprule
Model Feature & L1 & L2 & L3 \\
\midrule
Aggregated heat balance, Eq.~\eqref{eq:balance_L1}
  & \checkmark & -- & -- \\
Nodal heat balances, Eq.~\eqref{eq:balance_prod}--\eqref{eq:balance_consumer}
  & -- & \checkmark & \checkmark \\
Pipe thermal losses, Eq.~\eqref{eq:loss_supply}--\eqref{eq:loss_return}
  & -- & \checkmark & \checkmark \\
Global electricity balance, Eq.~\eqref{eq:balance_el}
  & \checkmark & \checkmark & \checkmark \\
Heat pump (dual-source), Eq.~\eqref{eq:hp_split}--\eqref{eq:hp_cap}
  & \checkmark & \checkmark & \checkmark \\
Thermal storage, Eq.~\eqref{eq:storage_soc}
  & \checkmark & \checkmark & \checkmark \\
Gas boiler \& electrode boiler
  & \checkmark & \checkmark & \checkmark \\
Grid buy/sell logic (Big-M), Eq.~\eqref{eq:grid_bigm}
  & \checkmark & \checkmark & \checkmark \\
Peak import demand charge
  & \checkmark & \checkmark & \checkmark \\
Hourly \ce{CO2} accounting
  & \checkmark & \checkmark & \checkmark \\
\midrule
\textit{Disabled (experimental control):} & & & \\
Transport delays / Pressure drop / Investment sizing
  & -- & -- & -- \\
\bottomrule
\end{tabular}
\end{table}

\subsection{MILP Formulation}
\label{sec:milp}

\subsubsection{Objective Function}

Total annual operational cost is minimized:
\begin{equation}
\min \; Z = C^{\mathrm{en}} + C^{\mathrm{fuel}}
+ C^{\mathrm{dump}} + C^{\ce{CO2}} + C^{\mathrm{dem}}
\label{eq:objective}
\end{equation}

\textbf{Electricity procurement and sales:}
\begin{equation}
C^{\mathrm{en}} = \Delta t \sum_{t \in \mathcal{T}} \left(
P_t^{\mathrm{buy}} \lambda_t^{\mathrm{buy}}
- P_t^{\mathrm{sell}} \lambda_t^{\mathrm{sell}} \right)
\label{eq:cost_energy}
\end{equation}
where $\lambda_t^{\mathrm{buy}} = \lambda_t^{\mathrm{spot}}
+ \lambda^{\mathrm{fee}} + \lambda^{\mathrm{grid}}$ includes spot
price, grid fees, and levies. The selling price accounts for bid-ask
spread, floor, haircut, and export fee (full derivation in
\Cref{app:selling_price}).

\textbf{Fuel cost:}
\begin{equation}
C^{\mathrm{fuel}} = \Delta t \sum_{g \in \mathcal{G}^{\mathrm{fuel}}}
\sum_{t \in \mathcal{T}} F_{g,t} \lambda_g^{\mathrm{fuel}}
\label{eq:cost_fuel}
\end{equation}

\textbf{Heat curtailment penalty:}
\begin{equation}
C^{\mathrm{dump}} = \Delta t \sum_{n \in \mathcal{N}}
\sum_{t \in \mathcal{T}} Q_{n,t}^{\mathrm{dump}} \lambda^{\mathrm{dump}}
\label{eq:cost_dump}
\end{equation}

\textbf{\ce{CO2} cost:}
\begin{equation}
C^{\ce{CO2}} = \frac{\lambda^{\ce{CO2}}}{1000} \left(
\sum_{t} E_t^{\ce{CO2},\mathrm{grid}}
+ \sum_{g,t} E_{g,t}^{\ce{CO2},\mathrm{fuel}} \right)
\label{eq:cost_co2}
\end{equation}

\textbf{Demand charge} (annual peak import fee, prorated):
\begin{equation}
C^{\mathrm{dem}} = \lambda^{\mathrm{dem}} \cdot f^{\mathrm{yr}}
\cdot P^{\mathrm{peak}}
\label{eq:cost_demand}
\end{equation}

\subsubsection{Energy Balance Constraints}
\label{sec:model_balance}

\textbf{L1 --- Aggregated copperplate balance (no losses):}
\begin{equation}
\sum_{g} Q_{g,t}^{\mathrm{th}}
+ \sum_{s} Q_{s,t}^{\mathrm{dis}}
= Q_t^{\mathrm{dem}} + Q_t^{\mathrm{dump}}
+ \sum_{s} Q_{s,t}^{\mathrm{ch}}
\qquad \forall t
\label{eq:balance_L1}
\end{equation}

\textbf{L2/L3 --- Nodal balance, producer nodes:}
\begin{equation}
\sum_{g \in \mathcal{G}_n} Q_{g,t}^{\mathrm{th}}
+ \sum_{s \in \mathcal{S}_n} Q_{s,t}^{\mathrm{dis}}
= Q_{n,t}^{\mathrm{dem}} + Q_{n,t}^{\mathrm{dump}}
+ \sum_{s \in \mathcal{S}_n} Q_{s,t}^{\mathrm{ch}}
+ Q_{n,t}^{\mathrm{loss}}
\quad \forall n \in \mathcal{N}^{\mathrm{prod}},\, \forall t
\label{eq:balance_prod}
\end{equation}

\textbf{L2/L3 --- Nodal balance, consumer nodes:}
\begin{equation}
Q_{n,t}^{\mathrm{pipe}}
+ \sum_{g \in \mathcal{G}_n} Q_{g,t}^{\mathrm{th}}
+ \sum_{s \in \mathcal{S}_n} Q_{s,t}^{\mathrm{dis}}
= Q_{n,t}^{\mathrm{dem}} + Q_{n,t}^{\mathrm{dump}}
+ \sum_{s \in \mathcal{S}_n} Q_{s,t}^{\mathrm{ch}}
\quad \forall n \in \mathcal{N}^{\mathrm{con}},\, \forall t
\label{eq:balance_consumer}
\end{equation}

The contrast between Eqs.~\eqref{eq:balance_L1} and
\eqref{eq:balance_prod} encodes the central experimental variable:
L1's missing loss term $Q_{n,t}^{\mathrm{loss}}$ is the sole
structural difference driving the L1$\rightarrow$L2 gap.

\textbf{Global electricity balance:}
\begin{equation}
P_t^{\mathrm{buy}} + \sum_{g \in \mathcal{G}^{\mathrm{CHP}}}
P_{g,t}^{\mathrm{el}}
= \sum_{h \in \mathcal{H}} P_{h,t}^{\mathrm{el}}
+ \sum_{r \in \mathcal{R}} P_{r,t}^{\mathrm{el}}
+ P_t^{\mathrm{sell}}
\qquad \forall t
\label{eq:balance_el}
\end{equation}

Grid import/export mutual exclusivity (Big-M):
\begin{equation}
P_t^{\mathrm{buy}} \leq M^{\mathrm{grid}} z_t, \qquad
P_t^{\mathrm{sell}} \leq M^{\mathrm{grid}} (1 - z_t)
\qquad \forall t
\label{eq:grid_bigm}
\end{equation}

\begin{equation}
P_t^{\mathrm{buy}} \leq \bar{P}^{\mathrm{imp}}, \qquad
P_t^{\mathrm{sell}} \leq \bar{P}^{\mathrm{exp}}
\qquad \forall t
\label{eq:grid_cap}
\end{equation}

Peak import tracking:
\begin{equation}
P^{\mathrm{peak}} \geq P_t^{\mathrm{buy}} \qquad \forall t
\label{eq:peak}
\end{equation}

\subsubsection{Thermal Loss Model (L2/L3 Only)}
\label{sec:model_losses}

Losses are computed separately for supply and return pipes using
steady-state conduction:
\begin{align}
Q_{p,t}^{\mathrm{loss,sup}} &= \frac{U_p L_p \big(
T_{\mathrm{sup}}^{\mathrm{nom}}(t) - T_{\mathrm{ground}}(t)\big)}{10^6}
\quad [\mathrm{MW}]
\label{eq:loss_supply} \\
Q_{p,t}^{\mathrm{loss,ret}} &= \frac{U_p L_p \big(
T_{\mathrm{ret}}^{\mathrm{nom}} - T_{\mathrm{ground}}(t)\big)}{10^6}
\quad [\mathrm{MW}]
\label{eq:loss_return}
\end{align}
where $U_p$~[\si{W/(m.K)}] is the linear heat transfer coefficient,
$L_p$~[m] the pipe length, and $T_{\mathrm{ground}}(t)$ the measured
ground temperature at installation depth. Losses are time-varying
through $T_{\mathrm{ground}}(t)$ and the heating-curve-dependent
$T_{\mathrm{sup}}^{\mathrm{nom}}(t)$, with seasonal amplitude of
approximately \placeholder{X}~MW\textsubscript{th}. For L1,
$Q_{p,t}^{\mathrm{loss}} = 0$ by design. The steady-state
approximation is valid for Peclet numbers $Pe \gg 1$; at case study
flow velocities (0.5--1.5~m/s), $Pe$ ranges from
\placeholder{XX--XX}~\cite{vanderHeijde2017}.

\subsubsection{Component Models}
\label{sec:model_components}

\textbf{Heat pumps.} Each heat pump $h$ supports dual-source
operation: waste heat recovery (WRG) with time-varying COP, and a
default source (ambient/ground) with constant COP:
\begin{align}
Q_{h,t} &= Q_{h,t}^{\mathrm{wrg}} + Q_{h,t}^{\mathrm{def}}
\label{eq:hp_split} \\
P_{h,t}^{\mathrm{el}} &= \frac{Q_{h,t}^{\mathrm{wrg}}}{
\mathrm{COP}_{h,t}^{\mathrm{wrg}}}
+ \frac{Q_{h,t}^{\mathrm{def}}}{\mathrm{COP}_h^{\mathrm{def}}}
\label{eq:hp_pel}
\end{align}
$\mathrm{COP}_{h,t}^{\mathrm{wrg}}$ is pre-computed via bilinear
interpolation of manufacturer lookup tables, clamped to $[0.5, 15.0]$
($<$0.3\% of timesteps affected). Interpolation error is bounded at
$\pm$2--3\% relative COP (\Cref{app:cop}), translating to $\pm$1--2\%
system cost---well below the topology gaps. Capacity and minimum
part-load:
\begin{equation}
\phi_h^{\mathrm{min}} \bar{Q}_h u_{h,t} \leq Q_{h,t}
\leq \bar{Q}_h u_{h,t}
\qquad \forall h,\, \forall t
\label{eq:hp_cap}
\end{equation}

\textbf{Thermal storage.} State-of-charge dynamics with
multiplicative self-discharge:
\begin{equation}
E_{s,t} = (1-\alpha_s)^{\Delta t} E_{s,t-1}
+ \eta_{s,t}^{\mathrm{ch}} Q_{s,t}^{\mathrm{ch}} \Delta t
- \frac{Q_{s,t}^{\mathrm{dis}} \Delta t}{\eta_{s,t}^{\mathrm{dis}}}
\qquad \forall s,\, \forall t
\label{eq:storage_soc}
\end{equation}
with energy capacity $E_{s,t} \leq \bar{E}_s a_{s,t}$, power
capacity $Q_{s,t}^{\mathrm{ch/dis}} \leq \bar{P}_s
\delta_{s,t}^{\mathrm{ch/dis}}$, and simultaneous charge/discharge
prevention $\delta_{s,t}^{\mathrm{ch}} + \delta_{s,t}^{\mathrm{dis}}
\leq a_{s,t}$.

\textbf{Gas boiler and electrode boiler:}
\begin{align}
Q_{g,t}^{\mathrm{th}} &= \eta_g^{\mathrm{th}} F_{g,t},
\quad Q_{g,t}^{\mathrm{th}} \leq \bar{Q}_g
& \forall g \in \mathcal{G}^{\mathrm{th}}
\label{eq:gen_thermal} \\
Q_{r,t} = \eta_{r,t} P_{r,t}^{\mathrm{el}},
\quad Q_{r,t} \leq \bar{Q}_r
& \forall r \in \mathcal{R}
\label{eq:p2h}
\end{align}

\textbf{CHP unit:}
$Q_{g,t}^{\mathrm{th}} = \eta_g^{\mathrm{th}} F_{g,t}$ and
$P_{g,t}^{\mathrm{el}} = \eta_g^{\mathrm{el}} F_{g,t}$ with
fixed heat-to-power ratio.

\subsubsection{Emissions Accounting}

Grid electricity and fuel emissions:
\begin{align}
E_t^{\ce{CO2},\mathrm{grid}} &= P_t^{\mathrm{buy}}
\cdot e_t^{\mathrm{grid}} \cdot \Delta t
\label{eq:em_grid} \\
E_{g,t}^{\ce{CO2},\mathrm{fuel}} &= F_{g,t} \cdot e_g^{\mathrm{fuel}}
\cdot \Delta t
\label{eq:em_fuel}
\end{align}
CHP emissions are allocated using the efficiency-weighted
method~\cite{Lund2014_4GDH}, consistent with Austrian national
reporting.

\subsection{Demand Heterogeneity Index}
\label{sec:model_hi}

The demand heterogeneity index HI quantifies spatial non-uniformity
of demand, serving as the key predictor for the L2$\rightarrow$L3 gap:
\begin{equation}
d_n = \frac{\sum_t Q_{n,t}^{\mathrm{dem}}}{\sum_n \sum_t
Q_{n,t}^{\mathrm{dem}}}, \qquad
\mathrm{HI} = \frac{N}{N-1} \sum_{n=1}^{N}
\left(d_n - \frac{1}{N}\right)^2
\label{eq:hi}
\end{equation}
HI~$=0$ indicates uniform demand; HI~$=1$ indicates complete
concentration at a single node. The normalization $N/(N\!-\!1)$
ensures comparability across different node counts.

\subsection{Implementation}
\label{sec:model_impl}

The model is implemented in Python~3.10 using
Pyomo~\cite{pyomo}. L1/L2/L3 instantiation from a single
parameterized codebase via YAML configuration files. Solver:
Gurobi~10.0, MIP gap~$\leq$~0.01\%, identical settings and random
seeds across all runs. Cost differences $\geq$0.1\% are thus
attributable to model structure, not solver tolerance.

%% ============================================================
%% 4. CASE STUDIES
%% ============================================================
\section{Case Studies}
\label{sec:case_studies}

\subsection{Primary Case: Stadtbach Industrial Heat Network}
\label{sec:case_stadtbach}

The primary validation case is an Austrian industrial heat network
providing process heat to a multi-building site with grid electricity
coupling. \Cref{tab:stadtbach} summarizes installed components and
system parameters.

\begin{table}[htbp]
\centering
\caption{Stadtbach network: installed components and key parameters.
Specific values anonymized per industrial partner agreement.}
\label{tab:stadtbach}
\begin{tabular}{@{}lll@{}}
\toprule
Component / Parameter & Value & Notes \\
\midrule
\multicolumn{3}{@{}l}{\textit{Demand}} \\
Peak / annual heat demand
  & \placeholder{8--15}~MW / \placeholder{25--50}~GWh
  & Process heat \\
\midrule
\multicolumn{3}{@{}l}{\textit{Network}} \\
Total pipe length (supply+return)
  & \placeholder{3--8}~km & \\
Nodes: L3 / L2
  & 30 / 5 & GIS / N-S-E-W-Central \\
Demand heterogeneity (HI)
  & \placeholder{0.XX} & Eq.~\eqref{eq:hi} \\
Supply / return temperature
  & \placeholder{XX} / \placeholder{XX}~\si{\celsius} & \\
\midrule
\multicolumn{3}{@{}l}{\textit{Generation}} \\
Heat pump (dual-source)
  & \placeholder{3--6}~MW\textsubscript{th}
  & COP 2.5--4.0 \\
Gas boiler
  & \placeholder{8--12}~MW\textsubscript{th}
  & Backup/peak \\
Electrode boiler (P2H)
  & \placeholder{1--3}~MW\textsubscript{th}
  & $\eta = 0.99$ \\
\midrule
\multicolumn{3}{@{}l}{\textit{Flexibility}} \\
Thermal storage
  & \placeholder{50--150}~MWh
  & Stratified \\
\midrule
\multicolumn{3}{@{}l}{\textit{Grid coupling}} \\
Import / export capacity
  & \placeholder{XX} / \placeholder{XX}~MW & \\
\bottomrule
\end{tabular}
\end{table}

\subsection{L1/L2/L3 Implementation and Controlled Matching}
\label{sec:case_levels}

\textbf{L3 (30 nodes):} Full physical topology from GIS data.
Individual pipe parameters. Serves as the validated reference.

\textbf{L2 (5 nodes):} Geographic aggregation into five demand zones.
Total $U \cdot L$ product preserved from L3:
$\sum_{p \in \mathcal{P}^{\mathrm{L2}}} U_p L_p =
\sum_{p \in \mathcal{P}^{\mathrm{L3}}} U_p L_p$.
Demand at each L2 node equals the sum of constituent L3 demands.

\textbf{L1 (1 node):} Single bus. Total demand equals the sum of
all 30 nodal demands. No losses.

All three levels share identical installed capacities, storage
parameters, boundary conditions (prices, temperatures, demands),
and solver settings.

\subsection{Input Data}
\label{sec:case_data}

\begin{itemize}[nosep]
    \item \textbf{Electricity prices}: EPEX Austria day-ahead spot,
    hourly, 2022/2023 (high volatility, negative-price episodes).
    \item \textbf{Temperatures}: Site-measured ambient and ground
    temperature, hourly, used for COP and losses simultaneously.
    \item \textbf{Heat demand}: Measured hourly at meter level;
    disaggregated to L3 nodes, aggregated to L2/L1.
    \item \textbf{Grid \ce{CO2} intensity}: ENTSO-E Transparency
    Platform, Austrian hourly emission factors.
    \item \textbf{Fuel prices}: Austrian wholesale gas (TTF-indexed),
    monthly averages.
\end{itemize}

\subsection{Validation Protocol}
\label{sec:case_validation}

L3 is validated against measured operational data \emph{prior to} the
controlled comparison. Validation targets: annual gas consumption
(MAPE~$<$5\%), annual electricity purchase (MAPE~$<$5\%), monthly HP
operating hours (MAPE~$<$10\%), annual storage full-cycle equivalent
(MAPE~$<$15\%), and peak heat demand (MAPE~$<$5\%). L1 and L2 are
evaluated \emph{relative to validated L3}---the question is not
whether they independently match reality, but how much they deviate
from the highest-fidelity model.

\subsection{Synthetic Parametric Network}
\label{sec:case_synthetic}

To establish generalizability beyond a single case, a synthetic
tree-topology network is parameterized by four independent properties
(\Cref{tab:synthetic}). A tree topology is used because the Stadtbach
network and the majority of small- to medium industrial DH installations
are radially structured. The full experiment comprises 36 feasible
configurations $\times$ 3 levels = 108 MILP solves. All configuration
files are published on Zenodo~\cite{Zenodo_data}.

\begin{table}[htbp]
\centering
\caption{Synthetic network parameter space. Ranges based on European
DH inventory~\cite{Werner2017,Frederiksen2013}. Of 81 combinations,
36 are retained after removing physically infeasible cases.}
\label{tab:synthetic}
\begin{tabular}{@{}llllp{4.0cm}@{}}
\toprule
Parameter & Low & Med & High & Rationale \\
\midrule
Pipe length [km]  & 1  & 5  & 15  & Campus to urban district \\
Demand HI          & 0.1 & 0.4 & 0.8 & Uniform to concentrated \\
Node count         & 5  & 15 & 30  & Sparse to dense \\
Storage ratio [h]  & 2  & 6  & 12  & Buffer to shift capacity \\
\bottomrule
\end{tabular}
\end{table}

%% ============================================================
%% 5. RESULTS
%% ============================================================
\section{Results}
\label{sec:results}

This section follows the causal chain of the topology effect:
validation confirms L3 as a reliable reference (\Cref{sec:results_validation}),
the L1$\rightarrow$L2 gap reveals the physics-inclusion effect
(\Cref{sec:results_cost}), storage dispatch provides the mechanistic
explanation (\Cref{sec:results_storage}), synthetic analysis establishes
generalizability (\Cref{sec:results_generalizability}), and sensitivity
analysis confirms robustness (\Cref{sec:results_sensitivity}).

\subsection{Model Validation (L3 vs.\ Measured Data)}
\label{sec:results_validation}

\begin{table}[htbp]
\centering
\caption{L3 model validation against measured operational data.}
\label{tab:validation}
\begin{tabular}{@{}lllll@{}}
\toprule
Variable & Measured & L3 Model & MAPE & Target \\
\midrule
Annual gas consumption [MWh]
  & \placeholder{XX,XXX} & \placeholder{XX,XXX}
  & \placeholder{X.X\%} & $<$5\% \\
Annual elec.\ purchase [MWh]
  & \placeholder{XX,XXX} & \placeholder{XX,XXX}
  & \placeholder{X.X\%} & $<$5\% \\
Monthly HP hours [h]
  & \placeholder{X,XXX} & \placeholder{X,XXX}
  & \placeholder{X.X\%} & $<$10\% \\
Storage full-cycle equiv.\ [--]
  & \placeholder{XXX} & \placeholder{XXX}
  & \placeholder{X.X\%} & $<$15\% \\
Peak heat demand [MW]
  & \placeholder{XX} & \placeholder{XX}
  & \placeholder{X.X\%} & $<$5\% \\
\bottomrule
\end{tabular}
\end{table}

\placeholder{[Validation commentary: all targets met, L3 accepted as
reference for controlled comparison.]}

\subsection{Cost and Emissions: L1 vs.\ L2 vs.\ L3}
\label{sec:results_cost}

\begin{table*}[htbp]
\centering
\caption{Annual operational cost and \ce{CO2} emissions across topology
levels. L3 is the validated reference. Negative $\Delta$ indicates
underestimation relative to L3.}
\label{tab:cost_co2}
\begin{tabular}{@{}llllllll@{}}
\toprule
Level & Total [\euro/a] & $\Delta$~Cost
  & Fuel [\euro/a] & Elec.\ [\euro/a] & \ce{CO2} [\euro/a]
  & Total \ce{CO2} [t/a] & $\Delta$~\ce{CO2} \\
\midrule
L1 & \placeholder{X,XXX,XXX} & $-$\placeholder{X.X}\%
  & \placeholder{XXX,XXX} & \placeholder{XXX,XXX}
  & \placeholder{XX,XXX}
  & \placeholder{X,XXX} & $-$\placeholder{X.X}\% \\
L2 & \placeholder{X,XXX,XXX} & $-$\placeholder{X.X}\%
  & \placeholder{XXX,XXX} & \placeholder{XXX,XXX}
  & \placeholder{XX,XXX}
  & \placeholder{X,XXX} & $-$\placeholder{X.X}\% \\
L3 & \placeholder{X,XXX,XXX} & (ref.)
  & \placeholder{XXX,XXX} & \placeholder{XXX,XXX}
  & \placeholder{XX,XXX}
  & \placeholder{X,XXX} & (ref.) \\
\bottomrule
\end{tabular}
\end{table*}

The L1$\rightarrow$L2 gap arises from a single mechanism: L1 does not
require generation to compensate network losses. At the Stadtbach
network's loss fraction of approximately \placeholder{X}\%, L1
systematically underestimates fuel consumption, electricity demand for
heat pumps, and peak grid import. The L2$\rightarrow$L3 gap arises
from spatially heterogeneous dispatch constraints that L2's five-zone
aggregation cannot resolve.

\ce{CO2} underestimation compounds through two channels: L1
over-dispatches the heat pump (low-emission source) because
loss-driven generation overhead is invisible, and L1 misses boiler
activation events at peak-loss hours.

\subsection{Dispatch Patterns and the Storage Diagnostic}
\label{sec:results_storage}

\placeholder{[Figure~1: Load duration curves by generation source per
topology level. Three panels: heat pump, gas boiler, electrode boiler.]}

Key dispatch differences relative to L3:
\begin{itemize}[nosep]
    \item Heat pump full-load hours: L1 overestimates by
    \placeholder{X}\%.
    \item Gas boiler activation hours: L1 underestimates by
    \placeholder{X}\%---loss-driven peak events absent.
    \item Electrode boiler utilization: stable across levels
    (primarily price-driven).
\end{itemize}

\textbf{Storage as diagnostic variable.} The temporal arbitrage value
of storage depends on holding costs. In L1, no physical holding cost
exists beyond self-discharge---charging responds purely to electricity
price differentials. In L2/L3, losses introduce an implicit
time-varying holding cost: heat stored at $t_1$ for discharge at $t_2$
must additionally compensate for the loss overhead at $t_2$. The
effective arbitrage spread becomes
$\lambda_{t_1}^{\mathrm{buy}} - (\lambda_{t_2}^{\mathrm{buy}} +
\lambda_{t_2}^{\mathrm{loss,marginal}})$, where
$\lambda_{t_2}^{\mathrm{loss,marginal}}$ is the marginal generation
cost of covering losses at $t_2$. During cold winter afternoons this
penalty is highest, discouraging the deep overnight charging that L1
schedules.

\begin{table}[htbp]
\centering
\caption{Storage dispatch metrics. Charging hour shift relative to
L3 (winter, October--March).}
\label{tab:storage}
\begin{tabular}{@{}lllll@{}}
\toprule
Metric & L1 & L2 & L3 & Unit \\
\midrule
Annual equiv.\ full cycles
  & \placeholder{XXX} & \placeholder{XXX}
  & \placeholder{XXX} & -- \\
Mean charging hour (winter)
  & \placeholder{XX:00} & \placeholder{XX:00}
  & \placeholder{XX:00} & h \\
Charging hour shift vs.\ L3
  & \placeholder{+X.X} & \placeholder{+X.X} & 0 & h \\
Peak shaving effectiveness
  & \placeholder{XX}\% & \placeholder{XX}\%
  & \placeholder{XX}\% & -- \\
Mean SOC at 06:00 (winter)
  & \placeholder{XX}\% & \placeholder{XX}\%
  & \placeholder{XX}\% & \% of $\bar{E}_s$ \\
\bottomrule
\end{tabular}
\end{table}

\placeholder{[Figure~2: Mean winter-week SOC profile, L1 vs.\ L3.
L1 charges later (overnight minimum price); L3 charges earlier
(anticipating loss-amplified afternoon discharge cost).]}

\subsection{Network Losses and Computational Performance}
\label{sec:results_computation}

\begin{table}[htbp]
\centering
\caption{Annual thermal losses and computational performance.}
\label{tab:losses_computation}
\begin{tabular}{@{}lllll@{}}
\toprule
Level & Annual loss [MWh] & Loss [\%] & Variables & Solve [s] \\
\midrule
L1 & 0 & 0 & \placeholder{X,XXX} & \placeholder{XX} \\
L2 & \placeholder{X,XXX} & \placeholder{X.X}
  & \placeholder{XX,XXX} & \placeholder{XXX} \\
L3 & \placeholder{X,XXX} & \placeholder{X.X}
  & \placeholder{XX,XXX} & \placeholder{X,XXX} \\
\bottomrule
\end{tabular}
\end{table}

\subsection{Generalizability: Synthetic Network Analysis}
\label{sec:results_generalizability}

\subsubsection{L1$\rightarrow$L2 Gap: Physics Inclusion}

The physics-inclusion gap is driven primarily by total pipe length.

\placeholder{[Figure~3a: L1$\rightarrow$L2 cost gap [\%] vs.\ total
pipe length [km], colored by storage ratio. Expected: near-linear
positive correlation.]}

For pipe lengths $<$1~km, annual losses are $<$1\% of demand and
the gap is negligible ($<$0.5\%). For 5--15~km, the gap grows to
\placeholder{X--Y}\%.

\subsubsection{L2$\rightarrow$L3 Gap: Spatial Refinement}

The spatial-refinement gap is driven by demand heterogeneity HI. For
HI~$<$~0.3, the gap remains below \placeholder{1}\% across all
configurations. For HI~$>$~0.5, it grows to \placeholder{X--Y}\%,
amplified by large storage (storage ratio~$>$~6~h), because incorrect
charging patterns compound over multiple daily cycles.

\placeholder{[Figure~3b: L2$\rightarrow$L3 gap [\%] as function of
HI (x-axis) and storage ratio (y-axis), medium pipe length.]}

\subsection{Sensitivity and Gap Stability}
\label{sec:results_sensitivity}

\placeholder{[Figure~4: Tornado diagram for L3 annual cost. Parameters:
gas price $\pm$20\%, electricity $\pm$20\%, CO$_2$ $\pm$50\%,
COP $\pm$10\%, self-discharge $\pm$50\%, ground temperature $\pm$3~K.]}

\begin{table}[htbp]
\centering
\caption{Stability of topology gaps across scenarios (Stadtbach,
L3 reference). The sign of the gap is stable in all scenarios.}
\label{tab:gap_stability}
\begin{tabular}{@{}lll@{}}
\toprule
Scenario & L1--L3 gap & L2--L3 gap \\
\midrule
Baseline         & \placeholder{X.X\%} & \placeholder{X.X\%} \\
High gas         & \placeholder{X.X\%} & \placeholder{X.X\%} \\
Low electricity  & \placeholder{X.X\%} & \placeholder{X.X\%} \\
High \ce{CO2}    & \placeholder{X.X\%} & \placeholder{X.X\%} \\
Cold year (+3~K) & \placeholder{X.X\%} & \placeholder{X.X\%} \\
Warm year ($-$3~K) & \placeholder{X.X\%} & \placeholder{X.X\%} \\
\bottomrule
\end{tabular}
\end{table}

The sign of both gaps is stable across all scenarios: L1 consistently
underestimates costs relative to L3, and the magnitude variation
remains within \placeholder{$\pm$X} percentage points of the baseline.

%% ============================================================
%% 6. DISCUSSION
%% ============================================================
\section{Discussion}
\label{sec:discussion}

\subsection{When Does Topology Detail Matter?}
\label{sec:disc_when}

The dominant finding is that the L1$\rightarrow$L2 transition drives
the majority of the topology effect. For systems with losses exceeding
3--5\% of annual heat supply---typical for Central European networks
with pipe lengths above 3~km---omitting losses produces systematic
cost underestimation of \placeholder{X--Y}\%. This is a direct
physical consequence: L1 requires less generation to meet the same
demand.

The L2$\rightarrow$L3 transition produces substantially smaller
differences for networks with uniform demand (HI~$<$~0.4). For the
large class of systems with moderate pipe length and near-uniform
demand, L2 provides a favorable fidelity--tractability trade-off.

\subsection{Actionable Topology Selection Criteria}
\label{sec:disc_criteria}

\Cref{tab:criteria} synthesizes findings from both the primary case
study and the synthetic analysis into directly applicable selection
rules.

\begin{table}[htbp]
\centering
\caption{Recommended topology selection criteria. Threshold:
topology gap $>$\placeholder{X}\% of L3 annual cost.}
\label{tab:criteria}
\begin{tabular}{@{}p{1.3cm}p{5.5cm}p{4.5cm}@{}}
\toprule
Level & Conditions & Rationale \\
\midrule
\textbf{L1}
  & Pipe $<$1~km \textbf{AND} no storage
  & Losses $<$1\%; no temporal arbitrage \\[4pt]
\textbf{L2}
  & Pipe 1--10~km \textbf{OR} storage~$\leq$6~h, \textbf{AND}
  HI~$\leq$~0.4
  & Losses significant; spatial distribution uniform \\[4pt]
\textbf{L3}
  & Pipe $>$10~km \textbf{OR} HI~$>$~0.5 \textbf{OR}
  storage $>$6~h with HI~$>$~0.3
  & Spatial dispatch differences material \\
\bottomrule
\end{tabular}
\end{table}

\textbf{Practical implication.} A practitioner using L1 for
operational planning would systematically schedule late overnight
storage charging. The physical network's higher afternoon losses
mean the true optimum requires earlier pre-charging. The economic
cost of this dispatch error is \placeholder{€X,XXX--XX,XXX}/year---a
recurring loss avoidable with an L2 model at a factor
\placeholder{X} increase in solve time.

\subsection{\ce{CO2} Policy Implications}
\label{sec:disc_co2}

Simplified topology models underestimate \ce{CO2} emissions through
two compounding channels: heat pump over-dispatch (low-emission source)
and missed boiler activations at peak-loss hours. A subtler third channel arises from temporal correlation:
loss-driven additional heat pump operation occurs at hours
whose grid emission factors may differ systematically from
the hours at which L1 schedules HP operation, introducing
a bias in hourly emission accounting. For industrial
networks subject to carbon reporting
obligations~\cite{EUheating2016}, L1-based emission accounting
introduces a systematic underestimation of
\placeholder{X--Y}~t\ce{CO2}/year. The present findings also
indicate that existing copperplate energy system frameworks (oemof,
PyPSA, FINE, Calliope) will exhibit this bias for DH systems with
significant pipe length.

\subsection{Limitations}
\label{sec:discussion_limitations}

\begin{itemize}[nosep]
    \item \textbf{Steady-state thermal model}: No transient startup
    dynamics; relevant for systems with frequent cycling.
    \item \textbf{No hydraulic coupling}: Pumping energy not optimized
    ($<$2\% of losses).
    \item \textbf{Hourly resolution}: Sub-hourly peaks not captured.
    \item \textbf{Single validation case}: Generalization relies on
    the synthetic parametric analysis.
    \item \textbf{Deterministic}: Uncertainty addressed via sensitivity
    analysis, not stochastic optimization.
    \item \textbf{Tree topology only}: Ring networks require additional
    flow direction binaries (deferred to future work).
    \item \textbf{Dispatch-only}: Potential topology effects on
    investment sizing not examined.
\end{itemize}

%% ============================================================
%% 7. CONCLUSION
%% ============================================================
\section{Conclusion}
\label{sec:conclusion}

This paper presents the first controlled comparison isolating topology
abstraction as the sole variable in district heating MILP dispatch
optimization. Three levels (L1 copperplate, L2 simplified multi-node,
L3 detailed multi-node) are embedded in a unified framework with
identical physics at L2/L3, validated against an Austrian industrial
heat network.

\textbf{Principal findings:}
\begin{enumerate}[nosep]
    \item The L1$\rightarrow$L2 transition dominates the topology
    effect: operational cost underestimation of \placeholder{X}\%,
    \ce{CO2} underestimation of \placeholder{X}\%. The mechanism is
    loss-driven generation underestimation.

    \item The L2$\rightarrow$L3 transition is marginal
    ($<$\placeholder{X}\%) for uniform demand (HI~$<$~0.4) but grows
    substantially for concentrated demand (HI~$>$~0.5).

    \item Storage charging hour distribution is the most sensitive
    dispatch diagnostic, shifted by \placeholder{X}~h between L1
    and L3 in winter due to loss-dependent holding costs.

    \item The L1$\rightarrow$L2 gap scales with pipe length
    ($R^2\!=\,$\placeholder{XX}); the L2$\rightarrow$L3 gap scales
    with demand heterogeneity---yielding actionable selection criteria
    (\Cref{tab:criteria}).
\end{enumerate}

\textbf{Recommendations:}
\begin{itemize}[nosep]
    \item L2 as minimum fidelity for any system with storage or pipe
    length $>$1~km.
    \item L3 when HI~$>$~0.4 or storage ratio~$>$~6~h with
    heterogeneous demand.
    \item Validate storage dispatch patterns as a topology adequacy
    diagnostic before reporting cost or emission results.
\end{itemize}

\textbf{Future work:} Ring topologies; stochastic variants;
investment--dispatch co-optimization; sub-hourly transient
validation.

%% ============================================================
%% REFERENCES
%% ============================================================
\bibliographystyle{elsarticle-num}
\begin{thebibliography}{99}

\bibitem{EUheating2016}
European Commission, An EU Strategy on Heating and Cooling,
COM(2016) 51 final, Brussels, 2016.

\bibitem{Lund2014_4GDH}
H.\ Lund, S.\ Werner, R.\ Wiltshire, S.\ Svendsen, J.E.\ Thorsen,
F.\ Hvelplund, B.V.\ Mathiesen,
4th Generation District Heating (4GDH),
Energy 68 (2014) 1--11.

\bibitem{Werner2017}
S.\ Werner,
District heating and cooling in Sweden,
Energy 126 (2017) 419--429.

\bibitem{Frederiksen2013}
S.\ Frederiksen, S.\ Werner,
District Heating and Cooling, Studentlitteratur, Lund, 2013.

\bibitem{Talebi2016}
B.\ Talebi, P.A.\ Mirzaei, A.\ Bastani, F.\ Haghighat,
A Review of District Heating Systems: Modeling and Optimization,
Frontiers in Built Environment 2 (2016) 22.

\bibitem{Buffa2019}
S.\ Buffa, M.\ Canova, M.\ D'Antoni, M.\ Baratieri, R.\ Fedrizzi,
5th generation district heating and cooling systems,
Renewable and Sustainable Energy Reviews 104 (2019) 504--522.

\bibitem{Lund2017_SmartEnergy}
H.\ Lund, P.A.\ Østergaard, D.\ Connolly, B.V.\ Mathiesen,
Smart energy and smart energy systems,
Energy 137 (2017) 556--565.

\bibitem{Fang2016}
T.\ Fang, R.\ Lahdelma,
Optimization of combined heat and power production with heat storage,
Applied Energy 162 (2016) 723--732.

\bibitem{Guelpa2019review}
E.\ Guelpa, V.\ Verda,
Thermal energy storage in district heating and cooling systems,
Applied Energy 252 (2019) 113474.

\bibitem{Hilpert2018_oemof}
S.\ Hilpert, C.\ Kaldemeyer, U.\ Krien, S.\ Günther,
C.\ Wingenbach, G.\ Plessmann,
The Open Energy Modelling Framework (oemof),
Energy Strategy Reviews 22 (2018) 16--25.

\bibitem{Krien2020_oemofthermal}
U.\ Krien, S.\ Günther, J.\ Launer, D.\ Möller,
oemof.thermal, GitHub repository, 2020.

\bibitem{Brown2018_pypsa}
T.\ Brown, J.\ Hörsch, D.\ Schlachtberger,
PyPSA: Python for Power System Analysis,
Journal of Open Research Software 6 (2018) 4.

\bibitem{Welder2018_FINE}
L.\ Welder, D.\ Ryberg, L.\ Kotzur, T.\ Grube, M.\ Robinius,
D.\ Stolten,
Spatio-temporal optimization of a future energy system,
Energy 158 (2018) 1130--1149.

\bibitem{Pfenninger2018_calliope}
S.\ Pfenninger, B.\ Pickering,
Calliope: a multi-scale energy systems modelling framework,
Journal of Open Source Software 3 (2018) 825.

\bibitem{Dorfner2016_urbs}
J.\ Dorfner,
urbs: A linear optimisation model for distributed energy systems,
GitHub repository, 2016.

\bibitem{Wiese2018_balmorel}
F.\ Wiese et al.,
Balmorel open source energy system model,
Energy Strategy Reviews 20 (2018) 26--34.

\bibitem{Connolly2014_EnergyPLAN}
D.\ Connolly, H.\ Lund, B.V.\ Mathiesen,
Smart Energy Europe,
Renewable and Sustainable Energy Reviews 60 (2016) 1634--1653.

\bibitem{Vesterlund2017}
M.\ Vesterlund, A.\ Toffolo, J.\ Dahl,
Optimization of multi-source complex district heating network,
Energy 126 (2017) 53--63.

\bibitem{Gabrielli2018}
P.\ Gabrielli, M.\ Gazzani, E.\ Martelli, M.\ Mazzotti,
Optimal design of multi-energy systems with seasonal storage,
Applied Energy 219 (2018) 408--424.

\bibitem{Wirtz2020}
M.\ Wirtz, L.\ Kivilip, P.\ Remmen, D.\ Müller,
5th Generation District Heating,
Applied Energy 260 (2020) 114158.

\bibitem{Mertz2016}
T.\ Mertz, S.\ Serra, A.\ Henon, J.C.\ Reneaume,
A MINLP optimization of district heating network design,
Energy 116 (2016) 1156--1172.

\bibitem{Buenning2020}
F.\ Bünning, M.\ Wetter, M.\ Fuchs, D.\ Müller,
Bidirectional low temperature district energy systems,
Applied Energy 209 (2018) 502--515.

\bibitem{Giraud2015}
L.\ Giraud, M.\ Merabet, R.\ Bavière, M.\ Vallée,
Optimal control of district heating systems,
Proceedings of the 12th International Modelica Conference, 2015.

\bibitem{Schweiger2017}
G.\ Schweiger, P.-O.\ Larsson, F.\ Magnusson, P.\ Lauenburg,
S.\ Velut,
District heating and cooling systems -- Framework for simulation
and dynamic optimization,
Energy 137 (2017) 566--578.

\bibitem{vanderHeijde2017}
B.\ van der Heijde et al.,
Dynamic equation-based thermo-hydraulic pipe model for district
heating and cooling systems,
Energy Conversion and Management 151 (2017) 158--169.

\bibitem{Blommaert2020}
M.\ Blommaert, J.\ Wack, M.\ Baelmans,
Adjoint optimization for topological design of district heating
networks,
Applied Energy 280 (2020) 116025.

\bibitem{Paalsson1999}
H.\ Pálsson, A.\ Larsen, B.N.\ Jørgensen, S.\ Bøhm, Y.\ Zhou,
Equivalent models of district heating systems,
Proceedings of the 7th International Symposium on District Heating
and Cooling, Lund, 1999.

\bibitem{Vandermeulen2018}
A.\ Vandermeulen, B.\ van der Heijde, L.\ Helsen,
Controlling district heating and cooling networks to unlock
flexibility,
Energy 151 (2018) 103--115.

\bibitem{Christidis2017}
A.\ Christidis, C.\ Koch, L.\ Pottel, G.\ Tsatsaronis,
The contribution of heat storage to profitable CHP operation,
Energy 41 (2012) 75--82.

\bibitem{pyomo}
W.E.\ Hart, C.\ Laird, J.-P.\ Watson, D.L.\ Woodruff,
Pyomo -- Optimization Modeling in Python, 2nd ed., Springer, 2017.

\bibitem{Paper2_CALION}
[Authors], CALION: A Flexible MILP Framework for Multi-Technology
District Heating Networks, [companion paper, in preparation, 2024].

\bibitem{Zenodo_data}
[Authors], Replication dataset: Topology abstraction effects on
district heating MILP dispatch, Zenodo [dataset], 2024.
\url{https://doi.org/10.5281/zenodo.XXXXXXX}

\end{thebibliography}

%% ============================================================
%% APPENDICES
%% ============================================================
\appendix

\section{Nomenclature}
\label{app:nomenclature}

\subsection*{Sets and Indices}
\begin{tabular}{@{}p{2.8cm}p{9.5cm}@{}}
$t \in \mathcal{T}$               & Time steps, $t = 1, \ldots, |\mathcal{T}|$ \\
$n \in \mathcal{N}$               & Network nodes \\
$\mathcal{N}^{\mathrm{prod/con}}$ & Producer / consumer nodes \\
$p \in \mathcal{P}$               & Pipe segments \\
$g \in \mathcal{G}$               & Thermal generators \\
$h \in \mathcal{H}$               & Heat pumps \\
$s \in \mathcal{S}$               & Thermal storage units \\
$r \in \mathcal{R}$               & Power-to-heat units \\
$f \in \mathcal{F}$  & Fuel types \\
$\mathcal{P}_n^{\mathrm{in/out}}$ & Pipes entering/leaving node $n$ \\
$\mathcal{G}_n,\, \mathcal{S}_n$  & Generators / storage at node $n$ \\
\end{tabular}

\subsection*{Parameters}
\begin{tabular}{@{}p{2.8cm}p{2.0cm}p{7.5cm}@{}}
$\Delta t$                              & h             & Time step duration \\
$U_p$                                   & W/(m$\cdot$K) & Pipe linear heat transfer coefficient \\
$L_p$                                   & m             & Pipe length \\
$T_{\mathrm{ground}}(t)$               & \si{\celsius} & Ground temperature \\
$T_{\mathrm{sup}}^{\mathrm{nom}}(t)$   & \si{\celsius} & Nominal supply temperature \\
$T_{\mathrm{ret}}^{\mathrm{nom}}$      & \si{\celsius} & Nominal return temperature \\
$\mathrm{COP}_{h,t}^{\mathrm{wrg}}$   & --            & WRG-source COP \\
$\mathrm{COP}_{h}^{\mathrm{def}}$      & --            & Default-source COP \\
$\eta_g^{\mathrm{th/el}}$              & --            & Thermal / electrical efficiency \\
$\eta_{s}^{\mathrm{ch/dis}}$           & --            & Storage charge / discharge efficiency \\
$\alpha_s$                             & h$^{-1}$      & Storage self-discharge rate \\
$a_{s,t}$                              & $\{0,1\}$     & Storage availability \\
$\lambda_t^{\mathrm{buy/sell}}$        & \euro/MWh     & Electricity buy / sell price \\
$\lambda_g^{\mathrm{fuel}}$            & \euro/MWh     & Fuel price \\
$\lambda^{\ce{CO2}}$                   & \euro/t       & Carbon price \\
$\lambda^{\mathrm{dem}}$               & \euro/MW/a    & Demand charge rate \\
$e_t^{\mathrm{grid}}$                  & kg/MWh        & Grid emission factor \\
$e_g^{\mathrm{fuel}}$                  & kg/MWh        & Fuel emission factor \\
$\bar{Q}_g,\, \bar{E}_s,\, \bar{P}_s$ & MW, MWh, MW   & Capacities \\
$\phi_h^{\mathrm{min}}$                & --            & HP minimum part-load ratio \\
$f^{\mathrm{yr}}$                      & --            & Year fraction \\
$\mathrm{HI}$                          & --            & Demand heterogeneity index \\
\end{tabular}

\subsection*{Decision Variables}
\begin{tabular}{@{}p{2.8cm}p{2.0cm}p{7.5cm}@{}}
$P_t^{\mathrm{buy/sell}}$         & MW        & Grid import / export \\
$P^{\mathrm{peak}}$               & MW        & Peak grid import \\
$Q_{g,t}^{\mathrm{th}}$          & MW        & Generator thermal output \\
$P_{h,t}^{\mathrm{el}}$          & MW        & HP electrical input \\
$Q_{h,t}^{\mathrm{wrg/def}}$     & MW        & HP output by source \\
$Q_{s,t}^{\mathrm{ch/dis}}$      & MW        & Storage charge / discharge \\
$E_{s,t}$                        & MWh       & Storage state of charge \\
$F_{g,t}$                        & MW        & Fuel input \\
$Q_{n,t}^{\mathrm{dump}}$        & MW        & Heat curtailment \\
$Q_{p,t}$                        & MW        & Pipe heat flow \\
$u_{h,t},\, z_t,\, \delta_{s,t}$ & $\{0,1\}$ & Binary: HP on/off, grid mode, storage mode \\
\end{tabular}

\section{Electricity Selling Price Formula}
\label{app:selling_price}

The complete electricity selling price, reflecting bilateral trading
conditions applicable to the Austrian industrial partner:
\begin{equation}
\lambda_t^{\mathrm{sell}} = \max\!\Bigg(\!
\underbrace{
\max\!\big(\lambda_t^{\mathrm{spot}} - \lambda^{\mathrm{spread}},\;
\lambda^{\mathrm{floor}}\big)
}_{\text{spread/floor-limited price}}
\cdot (1 - h^{\mathrm{sell}})
- \lambda^{\mathrm{sell,fee}} + \lambda^{\mathrm{sell,prem}},
\;\; 0\Bigg)
\end{equation}
$\lambda^{\mathrm{spread}}$: bid-ask spread;
$\lambda^{\mathrm{floor}}$: contractual minimum;
$h^{\mathrm{sell}} \in [0,1]$: haircut (transaction costs/imbalance
risk); $\lambda^{\mathrm{sell,fee}}$: export grid fee;
$\lambda^{\mathrm{sell,prem}}$: balancing premium. The outer
$\max(\cdot, 0)$ prevents selling penalties.

\section{COP Interpolation Error Analysis}
\label{app:cop}

Heat pump COP is pre-computed from a 2D manufacturer performance map
via bilinear interpolation ($\Delta T_s = \Delta T_a = 5$~K). The
interpolation error is bounded by:
\begin{equation}
|\mathrm{COP} - \widehat{\mathrm{COP}}| \leq
\frac{(\Delta T_s)^2}{8}
\max \left|\frac{\partial^2 \mathrm{COP}}{\partial T_s^2}\right|
+ \frac{(\Delta T_a)^2}{8}
\max \left|\frac{\partial^2 \mathrm{COP}}{\partial T_a^2}\right|
\end{equation}
Estimated from the data sheet: absolute COP error $\pm$0.08--0.12
($\approx$2--3\% at COP~$\approx$~3.5). Monte Carlo propagation
(1000 samples): 95\% CI of system cost $\pm$1.2\%, well below the
L1--L3 gap. Clamping to $[0.5, 15.0]$ activates in $<$0.3\% of
timesteps ($<$0.05\% annual cost effect).

\section{Anonymized Stadtbach Network Parameters}
\label{app:stadtbach}

\begin{table}[htbp]
\centering
\caption{Anonymized aggregate parameters, Stadtbach case study.}
\label{tab:app_params}
\begin{tabular}{@{}lll@{}}
\toprule
Parameter & Value & Unit \\
\midrule
Demand nodes (L3) / zones (L2)        & 30 / 5                    & -- \\
Total pipe length (supply+return)      & \placeholder{X.X}         & km \\
Pipe type distribution (DN100/200/250) & \placeholder{XX/XX/XX}    & \% \\
Mean $U$-value (supply / return)       & \placeholder{X.XX / X.XX} & W/(m$\cdot$K) \\
Supply / return temperature            & \placeholder{XX / XX}     & \si{\celsius} \\
Ground temp.\ (mean / range)           & \placeholder{X.X / X.X--X.X} & \si{\celsius} \\
Peak / annual demand                   & \placeholder{XX / XX}     & MW / GWh \\
Demand heterogeneity HI                & \placeholder{X.XX}        & -- \\
Storage $\bar{E}_s$ / ratio            & \placeholder{XXX / X.X}   & MWh / h \\
HP / gas boiler / P2H capacity         & \placeholder{X.X / XX / X.X} & MW\textsubscript{th} \\
Grid import / export capacity          & \placeholder{XX / XX}     & MW \\
\bottomrule
\end{tabular}
\end{table}

\end{document}